% Options for packages loaded elsewhere
% Options for packages loaded elsewhere
\PassOptionsToPackage{unicode}{hyperref}
\PassOptionsToPackage{hyphens}{url}
\PassOptionsToPackage{dvipsnames,svgnames,x11names}{xcolor}
%
\documentclass[
  letterpaper,
  DIV=11,
  numbers=noendperiod]{scrartcl}
\usepackage{xcolor}
\usepackage{amsmath,amssymb}
\setcounter{secnumdepth}{-\maxdimen} % remove section numbering
\usepackage{iftex}
\ifPDFTeX
  \usepackage[T1]{fontenc}
  \usepackage[utf8]{inputenc}
  \usepackage{textcomp} % provide euro and other symbols
\else % if luatex or xetex
  \usepackage{unicode-math} % this also loads fontspec
  \defaultfontfeatures{Scale=MatchLowercase}
  \defaultfontfeatures[\rmfamily]{Ligatures=TeX,Scale=1}
\fi
\usepackage{lmodern}
\ifPDFTeX\else
  % xetex/luatex font selection
\fi
% Use upquote if available, for straight quotes in verbatim environments
\IfFileExists{upquote.sty}{\usepackage{upquote}}{}
\IfFileExists{microtype.sty}{% use microtype if available
  \usepackage[]{microtype}
  \UseMicrotypeSet[protrusion]{basicmath} % disable protrusion for tt fonts
}{}
\makeatletter
\@ifundefined{KOMAClassName}{% if non-KOMA class
  \IfFileExists{parskip.sty}{%
    \usepackage{parskip}
  }{% else
    \setlength{\parindent}{0pt}
    \setlength{\parskip}{6pt plus 2pt minus 1pt}}
}{% if KOMA class
  \KOMAoptions{parskip=half}}
\makeatother
% Make \paragraph and \subparagraph free-standing
\makeatletter
\ifx\paragraph\undefined\else
  \let\oldparagraph\paragraph
  \renewcommand{\paragraph}{
    \@ifstar
      \xxxParagraphStar
      \xxxParagraphNoStar
  }
  \newcommand{\xxxParagraphStar}[1]{\oldparagraph*{#1}\mbox{}}
  \newcommand{\xxxParagraphNoStar}[1]{\oldparagraph{#1}\mbox{}}
\fi
\ifx\subparagraph\undefined\else
  \let\oldsubparagraph\subparagraph
  \renewcommand{\subparagraph}{
    \@ifstar
      \xxxSubParagraphStar
      \xxxSubParagraphNoStar
  }
  \newcommand{\xxxSubParagraphStar}[1]{\oldsubparagraph*{#1}\mbox{}}
  \newcommand{\xxxSubParagraphNoStar}[1]{\oldsubparagraph{#1}\mbox{}}
\fi
\makeatother


\usepackage{longtable,booktabs,array}
\newcounter{none} % for unnumbered tables
\usepackage{calc} % for calculating minipage widths
% Correct order of tables after \paragraph or \subparagraph
\usepackage{etoolbox}
\makeatletter
\patchcmd\longtable{\par}{\if@noskipsec\mbox{}\fi\par}{}{}
\makeatother
% Allow footnotes in longtable head/foot
\IfFileExists{footnotehyper.sty}{\usepackage{footnotehyper}}{\usepackage{footnote}}
\makesavenoteenv{longtable}
\usepackage{graphicx}
\makeatletter
\newsavebox\pandoc@box
\newcommand*\pandocbounded[1]{% scales image to fit in text height/width
  \sbox\pandoc@box{#1}%
  \Gscale@div\@tempa{\textheight}{\dimexpr\ht\pandoc@box+\dp\pandoc@box\relax}%
  \Gscale@div\@tempb{\linewidth}{\wd\pandoc@box}%
  \ifdim\@tempb\p@<\@tempa\p@\let\@tempa\@tempb\fi% select the smaller of both
  \ifdim\@tempa\p@<\p@\scalebox{\@tempa}{\usebox\pandoc@box}%
  \else\usebox{\pandoc@box}%
  \fi%
}
% Set default figure placement to htbp
\def\fps@figure{htbp}
\makeatother


% definitions for citeproc citations
\NewDocumentCommand\citeproctext{}{}
\NewDocumentCommand\citeproc{mm}{%
  \begingroup\def\citeproctext{#2}\cite{#1}\endgroup}
\makeatletter
 % allow citations to break across lines
 \let\@cite@ofmt\@firstofone
 % avoid brackets around text for \cite:
 \def\@biblabel#1{}
 \def\@cite#1#2{{#1\if@tempswa , #2\fi}}
\makeatother
\newlength{\cslhangindent}
\setlength{\cslhangindent}{1.5em}
\newlength{\csllabelwidth}
\setlength{\csllabelwidth}{3em}
\newenvironment{CSLReferences}[2] % #1 hanging-indent, #2 entry-spacing
 {\begin{list}{}{%
  \setlength{\itemindent}{0pt}
  \setlength{\leftmargin}{0pt}
  \setlength{\parsep}{0pt}
  % turn on hanging indent if param 1 is 1
  \ifodd #1
   \setlength{\leftmargin}{\cslhangindent}
   \setlength{\itemindent}{-1\cslhangindent}
  \fi
  % set entry spacing
  \setlength{\itemsep}{#2\baselineskip}}}
 {\end{list}}
\usepackage{calc}
\newcommand{\CSLBlock}[1]{\hfill\break\parbox[t]{\linewidth}{\strut\ignorespaces#1\strut}}
\newcommand{\CSLLeftMargin}[1]{\parbox[t]{\csllabelwidth}{\strut#1\strut}}
\newcommand{\CSLRightInline}[1]{\parbox[t]{\linewidth - \csllabelwidth}{\strut#1\strut}}
\newcommand{\CSLIndent}[1]{\hspace{\cslhangindent}#1}



\setlength{\emergencystretch}{3em} % prevent overfull lines

\providecommand{\tightlist}{%
  \setlength{\itemsep}{0pt}\setlength{\parskip}{0pt}}



 


\newcommand{\doctitle}{nietzsche: annotations}
\newcommand{\docdate}{2026-04-20}
\newcommand{\doclogoleft}{assets/fadenlogo-fwb.png}
\newcommand{\doclogoright}{assets/dcl-head-niemand.png}
% ============================================================
% loani.tex  —  include-in-header snippet, LuaLaTeX
% IPA handled by ipa.lua Lua filter at AST level.
% No ucharclasses, no babel, no runtime font switching.
% ============================================================

\usepackage{iftex}
\RequireLuaTeX

\usepackage{fontspec}

\setmainfont{Georgia}[
  Ligatures = TeX
]
% \usepackage[hebrew]{babel}
% Hebrew font — switched via ucharclasses (XeTeX) or
% directly referenced where needed
% \newfontfamily\hebrewfont{Alef}[
%   Script = Hebrew
% ]
% \newfontfamily\babelfont{Alef}[
%   Script = Hebrew
% ]
% \usepackage[hebrew, english]{babel}
% \babelfont[hebrew]{rm}{Alef}
% \usepackage[english, bidi=basic]{babel}
% \usepackage[english, hebrew, bidi=basic]{babel}
\usepackage[english, bidi=basic]{babel}

% \babelprovide[onchar=fonts ids, main=false]{hebrew} % all RTL, IPA wrongs
% \babelprovide[onchar=ids fonts, main=false]{english}
% \babelprovide[onchar=fonts ids, main=true]{english}
\babelprovide[onchar=fonts ids]{hebrew}

\babelfont[hebrew]{rm}{Alef}
% IPA font — called explicitly by ipa.lua filter output
\newfontfamily\ipafont{Source Sans 3}

% ------------------------------------------------------------
% Hebrew bidi — luatexja-style or direct LuaTeX primitive.
% Since babel is not loaded, we use the LuaTeX direction
% primitive directly: TRT = right-to-left, TLT = left-to-right
% Hebrew Unicode chars get RTL from the UBA via luaotfload.
% ------------------------------------------------------------
\directlua{
  luatexbase.add_to_callback("pre_linebreak_filter",
    function(head) return head end,
    "dummy_bidi")
}

\usepackage[hyperfootnotes=false]{hyperref}
% ============================================================
% header.tex  —  include-in-header snippet
% Replaces $var$ Pandoc syntax with \ifdefined\doc... guards.
% Vars injected via vars.lua Lua filter from YAML front matter:
%   \doctitle      \docsubtitle   \docauthor   \docdate
%   \docmeta       \docfoot       \doclogoleft \doclogoright
%
% No fancyhdr. Header fires via \AtBeginDocument,
% footer via \AtEndDocument — bidi-safe.
% ============================================================

\usepackage{graphicx}
\usepackage{tabularx}
\usepackage{xcolor}
\usepackage{array}
\usepackage{tikz}

% ---------- header font (adjust to taste) ----------
\newfontfamily\headerfont{Source Sans 3}

% ---------- footer colors ----------
\definecolor{footerbg}{HTML}{000000}
\definecolor{footerfg}{HTML}{FFFFFF}

% ---------- column types ----------
\newcolumntype{L}{>{\raggedright\arraybackslash}m{10mm}}
\newcolumntype{C}{>{\centering\arraybackslash}X}
\newcolumntype{R}{>{\raggedleft\arraybackslash}m{45mm}}

% vertical separator rule
\newcommand{\vsep}{\color{black}\vrule width 1pt}

% ============================================================
% CUSTOM HEADER
% ============================================================
\newcommand{\CustomHeader}{%
% \LTR{%
\noindent
\begin{tabularx}{\textwidth}{L !{\vsep} C !{\vsep} R}

% LEFT LOGO
\ifdefined\doclogoleft
  \includegraphics[width=10mm,height=10mm,keepaspectratio]{\doclogoleft}%
\fi
&
% CENTER: title + subtitle stacked in a minipage so \\ is a
% line break within the cell, not a new table row
\begin{minipage}[c]{\linewidth}
\centering
\ifdefined\doctitle
  {\Large\bfseries\headerfont\doctitle}%
\fi
\ifdefined\docsubtitle
  \\[2pt]{\normalsize\headerfont\docsubtitle}%
\fi
\vspace{5pt}

\end{minipage}
% &
% % CENTER TITLE BLOCK
% \ifdefined\doctitle
%   {\Large\bfseries\doctitle}%
% \fi
% \ifdefined\docsubtitle
%   \\\normalsize\docsubtitle
% \fi
&
% RIGHT LOGO
\ifdefined\doclogoright
  \includegraphics[height=10mm,keepaspectratio]{\doclogoright}%
\fi

\end{tabularx}

\hrule
% \vspace{6pt}

% ---------- META AREA ----------
{\small
\ifdefined\docauthor
  \textbf{Author:} \docauthor\\
\fi
\ifdefined\docdate
  \textbf{Date:} \docdate
\fi
\ifdefined\docmeta
  \\\docmeta\\
  snc: 16132.  
  {\ipafont ɪ ʃ ɔ ː ʁ}  
  {\ipafont wer es [pɛʟɛχɪʟ] ausspricht, und nicht [pɛʀɛχɪʟ]}
\fi
}
\vspace{6pt}

\hrule

% \vspace{12pt}
% }% end \LTR
}% end \CustomHeader

% ============================================================
% CUSTOM FOOTER  (overlay, anchored to page bottom)
% ============================================================
\newcommand{\CustomFooter}{%
\begin{tikzpicture}[remember picture, overlay]
\node[
  anchor=south,
  draw,
  fill=footerbg,
  text=footerfg,
  inner sep=10pt,
  rounded corners=2pt,
  text width=\dimexpr\textwidth-20pt\relax
]
at ([yshift=5mm]current page.south)
{%
  \begin{center}
  {\small
  \ifdefined\docfoot
    % {\docfoot}%
    {\headerfont\docfoot}%
  \fi
  }
  \end{center}
};
\end{tikzpicture}%
}% end \CustomFooter

% ============================================================
% FIRE header at doc start, footer at doc end
% ============================================================
% \AtBeginDocument{\CustomHeader}
\AtEndDocument{\CustomFooter}
\renewcommand{\maketitle}{}
\usepackage{booktabs}
\usepackage{longtable}
\usepackage{array}
\usepackage{multirow}
\usepackage{wrapfig}
\usepackage{float}
\usepackage{colortbl}
\usepackage{pdflscape}
\usepackage{tabu}
\usepackage{threeparttable}
\usepackage{threeparttablex}
\usepackage[normalem]{ulem}
\usepackage{makecell}
\usepackage{xcolor}
\KOMAoption{captions}{tableheading}
\makeatletter
\@ifpackageloaded{caption}{}{\usepackage{caption}}
\AtBeginDocument{%
\ifdefined\contentsname
  \renewcommand*\contentsname{Table of contents}
\else
  \newcommand\contentsname{Table of contents}
\fi
\ifdefined\listfigurename
  \renewcommand*\listfigurename{List of Figures}
\else
  \newcommand\listfigurename{List of Figures}
\fi
\ifdefined\listtablename
  \renewcommand*\listtablename{List of Tables}
\else
  \newcommand\listtablename{List of Tables}
\fi
\ifdefined\figurename
  \renewcommand*\figurename{Figure}
\else
  \newcommand\figurename{Figure}
\fi
\ifdefined\tablename
  \renewcommand*\tablename{Table}
\else
  \newcommand\tablename{Table}
\fi
}
\@ifpackageloaded{float}{}{\usepackage{float}}
\floatstyle{ruled}
\@ifundefined{c@chapter}{\newfloat{codelisting}{h}{lop}}{\newfloat{codelisting}{h}{lop}[chapter]}
\floatname{codelisting}{Listing}
\newcommand*\listoflistings{\listof{codelisting}{List of Listings}}
\makeatother
\makeatletter
\makeatother
\makeatletter
\@ifpackageloaded{caption}{}{\usepackage{caption}}
\@ifpackageloaded{subcaption}{}{\usepackage{subcaption}}
\makeatother
\usepackage{bookmark}
\IfFileExists{xurl.sty}{\usepackage{xurl}}{} % add URL line breaks if available
\urlstyle{same}
\hypersetup{
  colorlinks=true,
  linkcolor={blue},
  filecolor={Maroon},
  citecolor={cyan},
  urlcolor={Blue},
  pdfcreator={LaTeX via pandoc}}


\title{nietzsche: annotations}
\author{}
\date{2026-04-20}
\begin{document}
\maketitle

\CustomHeader


\subsection{notes on papers}\label{notes-on-papers}

\begin{verbatim}
Warning: 'xfun::attr()' is deprecated.
Use 'xfun::attr2()' instead.
See help("Deprecated")

Warning: 'xfun::attr()' is deprecated.
Use 'xfun::attr2()' instead.
See help("Deprecated")

Warning: 'xfun::attr()' is deprecated.
Use 'xfun::attr2()' instead.
See help("Deprecated")

Warning: 'xfun::attr()' is deprecated.
Use 'xfun::attr2()' instead.
See help("Deprecated")

Warning: 'xfun::attr()' is deprecated.
Use 'xfun::attr2()' instead.
See help("Deprecated")

Warning: 'xfun::attr()' is deprecated.
Use 'xfun::attr2()' instead.
See help("Deprecated")

Warning: 'xfun::attr()' is deprecated.
Use 'xfun::attr2()' instead.
See help("Deprecated")

Warning: 'xfun::attr()' is deprecated.
Use 'xfun::attr2()' instead.
See help("Deprecated")

Warning: 'xfun::attr()' is deprecated.
Use 'xfun::attr2()' instead.
See help("Deprecated")

Warning: 'xfun::attr()' is deprecated.
Use 'xfun::attr2()' instead.
See help("Deprecated")

Warning: 'xfun::attr()' is deprecated.
Use 'xfun::attr2()' instead.
See help("Deprecated")

Warning: 'xfun::attr()' is deprecated.
Use 'xfun::attr2()' instead.
See help("Deprecated")

Warning: 'xfun::attr()' is deprecated.
Use 'xfun::attr2()' instead.
See help("Deprecated")

Warning: 'xfun::attr()' is deprecated.
Use 'xfun::attr2()' instead.
See help("Deprecated")
\end{verbatim}

\paragraph{annotations: benn, dorische
welt.pdf}\label{annotations-benn-dorische-welt.pdf}

paper: Benn (1989)

\begin{verbatim}
Warning: 'xfun::attr()' is deprecated.
Use 'xfun::attr2()' instead.
See help("Deprecated")
\end{verbatim}

page

note

comment

1

Schlacht und Spitzkrüge, und Prinzchen, die phantastischen Kopfputz
tragen, doch ohne Blut und Jagd und ohne Roß und Waf-fen, an dies
Voreisenzeitalter im Tal von Knossos, diese un-geschützten Galerien,
illusionistisch aufgelöstenWände, zarten artistischen Stil, farbige
Fayencen, lange steife Röcke der Kreterinnen, enganliegende Taillen,
Busenhalter, femi-nine Treppen der Paläste mit niederen breiten Stufen,
be-quem für Weiberschritte -: grenzt über Mykene die dori-sche
Welt.(..)hethitischen Rassensplitter mit Mutterrecht, weib-

1

Dorische Welt Eine Untersuchung über die Beziehung von Kunst und Macht

\#cite\_benn\_gesammelte\_1989\_\_

1

Dorische Welt Eine Untersuchung über die Beziehung von Kunst und Macht

\#cite\_benn\_gesammelte\_1989\_\_

2

ihm was beibringen, ihm plausibel macnen.(..)Prunkreden, Reklamereden,
Reden gegen Entrée. Dann kam »das Bezaubern« hinzu, symmetrischer Bau
der Sätze, oleicheknende heinghe reimende Worte, ähnlich auslau-

2

ihm was beibringen, ihm plausibel macnen.(..)Prunkreden, Reklamereden,
Reden gegen Entrée. Dann kam »das Bezaubern« hinzu, symmetrischer Bau
der Sätze, oleicheknende heinghe reimende Worte, ähnlich auslau-

3

underte en оп, b pielerische und senone 5 gestellt, vert , auch ''
sische Welt: die Glieder des menschliche

3

underte en оп, b pielerische und senone 5 gestellt, vert , auch ''
sische Welt: die Glieder des menschliche

6

man es auch hinsichtlich staats-feindlicher Menschen machen, schrieb
Demokrit, und Pro-tagoras läßt den Gottvater sprechen: Gib in meinem
Na-men das Gesetz, daß man einen Menschen ohne sittliches Bewußtsein und
ohne Rechtsgefühl als cinen Krebsschaden des Gemeinwesens vernichten
soll -: hieraus entstand das 5. Jahrhundert, der größte Glanz der weißen
Rasse, der prä-gende, der absolute, nicht nur der mittelmeerisch
be-grenzte, in diesen 50 Jahren nach Salamis, alles gruppiert sich um
diese Seeschlacht: Aschylos kä

6

man es auch hinsichtlich staats-feindlicher Menschen machen, schrieb
Demokrit, und Pro-tagoras läßt den Gottvater sprechen: Gib in meinem
Na-men das Gesetz, daß man einen Menschen ohne sittliches Bewußtsein und
ohne Rechtsgefühl als cinen Krebsschaden des Gemeinwesens vernichten
soll -: hieraus entstand das 5. Jahrhundert, der größte Glanz der weißen
Rasse, der prä-gende, der absolute, nicht nur der mittelmeerisch
be-grenzte, in diesen 50 Jahren nach Salamis, alles gruppiert sich um
diese Seeschlacht: Aschylos kä

8

Körper zur Zucht: das Gesetz bestimmte das heirats-fahige Alter und
wählte den günstigsten Zeitpunkt und die günstigsten Umstände für eine
Schwängerung aus. Man ging wie in Gestüten vor, man vernichtete die
schlechtge-lungene Frucht. Der Körper zum Krieg, der Körper zum Fest,
der Körper zum Laster und der Körper endlich dann zur Kunst, das war die
dorische Saat und die hellenische Geschichte Dori

8

Dorische Welt war die größte griechisch tike Sittlichkeit, also siegende
Ordnung und von Göttern stammende Macht. Nichts wissen ihre Sagen von
Horten, Schätzen, hohlen Bergen, ihre Begehrlichkeit geht nicht auf
Gold, sondern auf heilige Dinge, magische Waffen aus Hephaistos' Hand,
auf das Goldene Vlies, das Halsband der Harmonia, das Zepter des Zeus.
Sparta, das war auch ein schick. Menschen, welche bei den übri-

8

inken der Turnzucht?«Das Sinken der Turnzucht - mit ihr sanken die
dorische Welt, Olympia, die graue Säule ohne Fuß und die der
Her-renschicht günstigen Orakel. Euripides ist skeptisch, cin-sam und
atheistisch, es steigen bei ihm bereits die Allge-meinbegriffe isoliert
auf: das Gute«, das Rechte«, »die Tugenda, die Bildung«; er ist
pazifistisch und antihero-isch: vor allem Frieden und keine sizilische
Expedition, er 296

8

Dorische Welt war die größte griechisch tike Sittlichkeit, also siegende
Ordnung und von Göttern stammende Macht. Nichts wissen ihre Sagen von
Horten, Schätzen, hohlen Bergen, ihre Begehrlichkeit geht nicht auf
Gold, sondern auf heilige Dinge, magische Waffen aus Hephaistos' Hand,
auf das Goldene Vlies, das Halsband der Harmonia, das Zepter des Zeus.
Sparta, das war auch ein schick. Menschen, welche bei den übri-

8

Körper zur Zucht: das Gesetz bestimmte das heirats-fahige Alter und
wählte den günstigsten Zeitpunkt und die günstigsten Umstände für eine
Schwängerung aus. Man ging wie in Gestüten vor, man vernichtete die
schlechtge-lungene Frucht. Der Körper zum Krieg, der Körper zum Fest,
der Körper zum Laster und der Körper endlich dann zur Kunst, das war die
dorische Saat und die hellenische Geschichte Dori

8

inken der Turnzucht?«Das Sinken der Turnzucht - mit ihr sanken die
dorische Welt, Olympia, die graue Säule ohne Fuß und die der
Her-renschicht günstigen Orakel. Euripides ist skeptisch, cin-sam und
atheistisch, es steigen bei ihm bereits die Allge-meinbegriffe isoliert
auf: das Gute«, das Rechte«, »die Tugenda, die Bildung«; er ist
pazifistisch und antihero-isch: vor allem Frieden und keine sizilische
Expedition, er 296

9

ld ihre Kampfspiele halten, sondern um die Trefflichkeit. « Diese
Trefflichkeit, dieser Kranz, dieses Kampffest zwischen den großen
Schlachten, das war hinter der panhellenischen Silhouctte die dorische
Welt.(..)

9

ld ihre Kampfspiele halten, sondern um die Trefflichkeit. « Diese
Trefflichkeit, dieser Kranz, dieses Kampffest zwischen den großen
Schlachten, das war hinter der panhellenischen Silhouctte die dorische
Welt.(..)

10

er cinzelnen dorischen Landschaften gesammelt und nach Kunstregeln
geordnet, hier wurde ein allgültiges Tonsystem festgelegt Hier gab es
von 645 an einen Gesetzgeber der Musik, Musik war gesetzliches Lehrfach
und mußte von allen Bewohnern bis zum 30. Lebensjahr betrieben werden,
alle mußten die Flöte spielen, durch Poesie und Gesang wurden die
Gesetze den Nachkommen überliefert, und man zog zum Kampf aus unter dem
Klang von Flöten, Lyren und Kitharen, Alkman treffliches Kitharspiel
geht dem Schwert voran«, sang Und

10

nen nicht zu trennen, das war die spartanische Sendung, und die
lakedä-monischen Gebräuche verdrängten die homerischen. Bald gibt es
keine Stadt mehr ohne ein Gymnasium, es ist eines der Zeichen, an denen
man eine griechische Stadt erkennt.(..)Aus cinem solchen Viereck mit
Säulenhallen und Platanen-alleen, gewöhnlich neben einer Quelle oder
neben cinem Bach, ging auch dic Akademie hervor, und die große
Philo-sophie entstand darin. In der späten Griechenzeit hicß es sogar,
die Spartaner hätten die gesunkene griechische Mu-sik dreimal gerettet,
und man stellte Sparta allegorisch als Weib mit einer Lyra dar. Es ist
auch Sparta, in der das erste Gebäude für musikalische und

10

denfalls.(..)wirklich eine Beziehung, die geschichtlich 400 Jahre dau
erte, von der öffentlichen Gewalt ganz unmittelbar zur Kunst, vom
Heroentum der äußeren Haltung und der Tat, von den Schlachtfeldern von
Marathon und Salamis zur Formfindung des letzten Parthenonstils, hier
kann man wirklich von einer Geburt der Kunst aus der Macht spre-chen, in
der Geschichte der Statue und an dieser Stelle je-es also war Sparta, so
sehr war es der Ausgangspunkt

10

nen nicht zu trennen, das war die spartanische Sendung, und die
lakedä-monischen Gebräuche verdrängten die homerischen. Bald gibt es
keine Stadt mehr ohne ein Gymnasium, es ist eines der Zeichen, an denen
man eine griechische Stadt erkennt.(..)Aus cinem solchen Viereck mit
Säulenhallen und Platanen-alleen, gewöhnlich neben einer Quelle oder
neben cinem Bach, ging auch dic Akademie hervor, und die große
Philo-sophie entstand darin. In der späten Griechenzeit hicß es sogar,
die Spartaner hätten die gesunkene griechische Mu-sik dreimal gerettet,
und man stellte Sparta allegorisch als Weib mit einer Lyra dar. Es ist
auch Sparta, in der das erste Gebäude für musikalische und

10

denfalls.(..)wirklich eine Beziehung, die geschichtlich 400 Jahre dau
erte, von der öffentlichen Gewalt ganz unmittelbar zur Kunst, vom
Heroentum der äußeren Haltung und der Tat, von den Schlachtfeldern von
Marathon und Salamis zur Formfindung des letzten Parthenonstils, hier
kann man wirklich von einer Geburt der Kunst aus der Macht spre-chen, in
der Geschichte der Statue und an dieser Stelle je-es also war Sparta, so
sehr war es der Ausgangspunkt

10

er cinzelnen dorischen Landschaften gesammelt und nach Kunstregeln
geordnet, hier wurde ein allgültiges Tonsystem festgelegt Hier gab es
von 645 an einen Gesetzgeber der Musik, Musik war gesetzliches Lehrfach
und mußte von allen Bewohnern bis zum 30. Lebensjahr betrieben werden,
alle mußten die Flöte spielen, durch Poesie und Gesang wurden die
Gesetze den Nachkommen überliefert, und man zog zum Kampf aus unter dem
Klang von Flöten, Lyren und Kitharen, Alkman treffliches Kitharspiel
geht dem Schwert voran«, sang Und

11

Sparta, die Macht.(..)Wir leiten' also aus Sparta Griechenland ab, und
aus dem Dorisch-Apollinischen die griechische Welt. Dionysos steht hier
wieder in den Grenzen, in denen er vor 1871 (Geburt der Tragödie aus der
Musik «) stand. Die Griechen waren ein primitives, d.~h. ein rauschnahes
Volk, ihr Zeus-dienst hatte orgiastische Züge; große, rauschhafte
Erre-gungswellen traten periodisch bei ihnen auf, auch in Sparta, viel
Kathartisches hatten sie aus den kretischen Schulen übernommen. Aber wir
haben inzwischen primitive Völ-ker aus Reisesc

11

sie auch immer wieder »der Erzieher« genannt, und sehr interessant ist,
wie Plato, geistig der letzte Dorer, der wäh-rend der Auflösung noch
einmal den Kampfgegen den In-dividualismus, das Schwermütig-Reizvolle
der Kunst, die süBliche Muße«, die Schattierkunst« aufnimmt, um für
mittlere Leben«, die »Gemeinde und »die vernünftigen Gedanken«, »das die
»Stadt mit der untadeligen Verfas-sung« zu kämpfen, diese Wehmut nach
Sparta ausdrückt.(..)im Theages« sagt er von einem tugendhaften Mann,
der über die Tugend Vorträge hält: »In der wunderbaren Har-monie seiner
Handlungen und seiner Worte erkennt man die dorische Weise wieder, die
einzige, welche wirklich griechisch ist. « Das war 500 Jahre nach
Lykurg. Und von hier aus versteht man nun auch seine w

11

chischen Welt entstand. Zwischen Rausch und Kunst m Apollo, die große
züchtende Kraft. Und da wir heute nicht mehr so wagnerisch erregt sind,
um das Bedürfnis zu ha-ben, Tristan schon in Thrazien nachzuweisen,
blicken wir nach Dorien, nicht nach Dion, fragen wir uns nach der grie-

11

chischen Welt entstand. Zwischen Rausch und Kunst m Apollo, die große
züchtende Kraft. Und da wir heute nicht mehr so wagnerisch erregt sind,
um das Bedürfnis zu ha-ben, Tristan schon in Thrazien nachzuweisen,
blicken wir nach Dorien, nicht nach Dion, fragen wir uns nach der grie-

11

Sparta, die Macht.(..)Wir leiten' also aus Sparta Griechenland ab, und
aus dem Dorisch-Apollinischen die griechische Welt. Dionysos steht hier
wieder in den Grenzen, in denen er vor 1871 (Geburt der Tragödie aus der
Musik «) stand. Die Griechen waren ein primitives, d.~h. ein rauschnahes
Volk, ihr Zeus-dienst hatte orgiastische Züge; große, rauschhafte
Erre-gungswellen traten periodisch bei ihnen auf, auch in Sparta, viel
Kathartisches hatten sie aus den kretischen Schulen übernommen. Aber wir
haben inzwischen primitive Völ-ker aus Reisesc

11

sie auch immer wieder »der Erzieher« genannt, und sehr interessant ist,
wie Plato, geistig der letzte Dorer, der wäh-rend der Auflösung noch
einmal den Kampfgegen den In-dividualismus, das Schwermütig-Reizvolle
der Kunst, die süBliche Muße«, die Schattierkunst« aufnimmt, um für
mittlere Leben«, die »Gemeinde und »die vernünftigen Gedanken«, »das die
»Stadt mit der untadeligen Verfas-sung« zu kämpfen, diese Wehmut nach
Sparta ausdrückt.(..)im Theages« sagt er von einem tugendhaften Mann,
der über die Tugend Vorträge hält: »In der wunderbaren Har-monie seiner
Handlungen und seiner Worte erkennt man die dorische Weise wieder, die
einzige, welche wirklich griechisch ist. « Das war 500 Jahre nach
Lykurg. Und von hier aus versteht man nun auch seine w

12

n.(..)Betrachtet man dies alles nicht unter moralischen, senti-mentalen,
geschichtsphilosophischen Gesichtspunkten, sondern nach
anthropologischen Grundsätzen, so sehen wir auf der einen Scite die
Macht und neben ihr den andern Ausbruch der hellenischen Volkheit: die
Kunst. Wie verhal-ten sich diese beiden nun zueinander, welches waren
ihre

12

n.(..)Betrachtet man dies alles nicht unter moralischen, senti-mentalen,
geschichtsphilosophischen Gesichtspunkten, sondern nach
anthropologischen Grundsätzen, so sehen wir auf der einen Scite die
Macht und neben ihr den andern Ausbruch der hellenischen Volkheit: die
Kunst. Wie verhal-ten sich diese beiden nun zueinander, welches waren
ihre

13

den, Fortschrittsinstrumentarium. Der Gedanke ist immer der Abkömmling
der Not«, sagt Schiller, bei dem wir ja ein sehr bewußtes Umlegen der
Achse vom moralischen zum ästhetischen Weltbild wahrnehmen, er meint,
der Gedanke steht immer nahe bei den Zweckmäßigkeiten und der Trieb- bei
Äxten und Morgenstern, er ist Natur, aber

13

Und dabei Form nie als Ermüd deutsch-bürgerlichen Sinne,sondern als die
enorme menschliche Macht, die Macht schlechthin, der Sieg über nackten
Tatbestand und zivilisatorische Sachverhalte, eben als das
Abendländische, die Uberhöhung, der reale eigen-kategoriale Geist, der
Ausgleich und die Sammlung der Fragmente. Nietzsche als Ganzes in einem
einzigen Satz, das könnte nur sein tiefster und zukünftigster sein: Nur
als ästhetisches Phänomen ist das Dasein und die Welt ewig
gerechtfertigt. « Das aber ist hellenisch.(..)uch Goethe sehen wir hier
stehen. Seine Iphigenie ist

13

den, Fortschrittsinstrumentarium. Der Gedanke ist immer der Abkömmling
der Not«, sagt Schiller, bei dem wir ja ein sehr bewußtes Umlegen der
Achse vom moralischen zum ästhetischen Weltbild wahrnehmen, er meint,
der Gedanke steht immer nahe bei den Zweckmäßigkeiten und der Trieb- bei
Äxten und Morgenstern, er ist Natur, aber

13

Und dabei Form nie als Ermüd deutsch-bürgerlichen Sinne,sondern als die
enorme menschliche Macht, die Macht schlechthin, der Sieg über nackten
Tatbestand und zivilisatorische Sachverhalte, eben als das
Abendländische, die Uberhöhung, der reale eigen-kategoriale Geist, der
Ausgleich und die Sammlung der Fragmente. Nietzsche als Ganzes in einem
einzigen Satz, das könnte nur sein tiefster und zukünftigster sein: Nur
als ästhetisches Phänomen ist das Dasein und die Welt ewig
gerechtfertigt. « Das aber ist hellenisch.(..)uch Goethe sehen wir hier
stehen. Seine Iphigenie ist

13

wenn wir uns jetzt einmal dem Wesen der griechischen Kunst zuwenden, so
drückt der dorische Tempel ja nichts aus, er ist nicht verständlich, und
die Säule ist nicht natür-lich, sie nehmen nicht einen konkreten
politischen oder kultischen Willen in sich auf, sie sind überhaupt mit
nichts parallel, sondern das Ganze ist ein S til, das heißt, es ist von
innen gesehen ein bestimmtes Raumgefühl, eine bestimmte Raumpanik, und
von außen gesehen sind es bestimmte An-lagen und Prinzipien, um das
darzustellen, es auszudrük-ken, also es zu beschwören. Dieses
Darstellungsprinzip stammt nicht mehr unmittelba

13

wenn wir uns jetzt einmal dem Wesen der griechischen Kunst zuwenden, so
drückt der dorische Tempel ja nichts aus, er ist nicht verständlich, und
die Säule ist nicht natür-lich, sie nehmen nicht einen konkreten
politischen oder kultischen Willen in sich auf, sie sind überhaupt mit
nichts parallel, sondern das Ganze ist ein S til, das heißt, es ist von
innen gesehen ein bestimmtes Raumgefühl, eine bestimmte Raumpanik, und
von außen gesehen sind es bestimmte An-lagen und Prinzipien, um das
darzustellen, es auszudrük-ken, also es zu beschwören. Dieses
Darstellungsprinzip stammt nicht mehr unmittelba

14

Alle Lust will Ewigkeit, sagte das vorige Jahrhundert, das neue fhrt
fort: Alle Ewigkeit will Kunst. Die absolute Kunst, die Form. Doch
palles Schöne ist schwer, und wer

14

Alle Lust will Ewigkeit, sagte das vorige Jahrhundert, das neue fhrt
fort: Alle Ewigkeit will Kunst. Die absolute Kunst, die Form. Doch
palles Schöne ist schwer, und wer

\paragraph{annotations: benn, kilpper 1943
brief.pdf}\label{annotations-benn-kilpper-1943-brief.pdf}

paper: Benn (1971)

\begin{verbatim}
Warning: 'xfun::attr()' is deprecated.
Use 'xfun::attr2()' instead.
See help("Deprecated")
\end{verbatim}

page

note

comment

1

AN G U S TAV K I L P P E R Landsberg

\#cite\_benn\_lyrik\_1971\_\_

1

AN G U S TAV K I L P P E R Landsberg

\#cite\_benn\_lyrik\_1971\_\_

\paragraph{annotations: Gesammelte Werke Bd. 3 Gedichte by Gottfried
Benn.pdf}\label{annotations-gesammelte-werke-bd.-3-gedichte-by-gottfried-benn.pdf}

paper: Benn (1960)

\begin{verbatim}
Warning: 'xfun::attr()' is deprecated.
Use 'xfun::attr2()' instead.
See help("Deprecated")
\end{verbatim}

page

note

comment

219

VERLORENES ICH

1948(..)rönne: 1916(..)(..)\#cite\_benn\_gesammelte\_1960\_\_(..)

219

auf deinem grauen Stein von Notre-Dame

wasserspeier notre dame

219

Katalaunen

neologism

219

Kaldaunen

kutteln, innereien:rindermagen

219

Lamm

lamda

219

Katalaunen

neologism

219

Lamm

lamda

219

Kaldaunen

kutteln, innereien:rindermagen

220

Woher, wohin - nicht Nacht, nicht Morgen, kein Evoe, kein Requiem, du
mochtest dir ein Stichwort borgen allein bei wem?

evoe: freudenschrei des dionyskultes

220

o feme zwingende erfiillte Stunde, die einst auch das verlorne Ich
umschlofi

marginnote4app://note/D263ECE4-4818-4065-B55B-16CAB9502871

220

als sich alle einer Mitte neigten

zerteilung, multiperspektive,individualisierung(..)1943: auch
gleichschaltung, NS, massen / sehnsucht nach führung; benn: kritik

220

als sich alle einer Mitte neigten

zerteilung, multiperspektive,individualisierung(..)1943: auch
gleichschaltung, NS, massen / sehnsucht nach führung; benn: kritik

220

o feme zwingende erfiillte Stunde, die einst auch das verlorne Ich
umschlofi

3

1

2

307

Von Blumen mufit du solche wahlen, die bluhn am Zaun und halb 1m Acker
schon, die in das Zimmer tun, die Laute zahlen des Lebens Laute, seinen
Ton:

307

Tiere, die Perlen bilden, sind versdilossen, sie liegen still und kennen
nur die See

308

der Aon traumt, der Aon ist ein Knabe, der mit sich selbst auf einem
Brette spielt

346

Ob Rosen, ob Schnee, ob Meere, was alles erbliihte, verblich, es gibt
nur zwei Dinge: die Leere und das gezeichnete Ich.

marginnote4app://note/80AF25DE-102C-4B8C-A51D-2E10DEED1BD5

\paragraph{annotations:
gesammelteschrif03wagn.pdf}\label{annotations-gesammelteschrif03wagn.pdf}

paper: @

\begin{verbatim}
Warning: 'xfun::attr()' is deprecated.
Use 'xfun::attr2()' instead.
See help("Deprecated")
\end{verbatim}

page

note

comment

20

Die Kunst und die Revolution

\#init\_Wagner, Die Kunst und die Revolution\_

20

Die Kunst und die Revolution

\#init\_Wagner, Die Kunst und die Revolution\_

21

2)er griecîfdê (SJeift, mie er fic§ gu feiner \^{}lütleêit in
(\textbf{taat?}) unb ûnft gu erfennen gab, fanb, nadfjbem er bie roi§e
\textsuperscript{}laturreligion ber afiatifcên §eimat§ übermunben, unb
hm \^{}'ömn unb ftarlea freien 3)Zenfdên auf bie ©piê feine§ religiöfen
êmujßtfeinS ge= fteHt Iâtte, feinen entfpredênbften 2lu§brudf in
SlpoUon, bem eigent\^{} liefen §aupt= unb 5flationaIgotte ber êHenifdên
©tämme.

\paragraph{annotations: Günther 2008, Der
Wettkampf.pdf}\label{annotations-gunther-2008-der-wettkampf.pdf}

paper: Günther (2008)

\begin{verbatim}
Warning: 'xfun::attr()' is deprecated.
Use 'xfun::attr2()' instead.
See help("Deprecated")
\end{verbatim}

page

note

comment

1

Günther

1

Friederike Felicitas Günther Rhythmus beim frühen Nietzsche

\#cite\_gunther\_rhythmus\_2008\_\_

5

Der Rhythmik gesteht er hier die Macht zu, den „Willen``, d.h. das
unablässige ewige dionysische Werden als Naturkraft ohne Grenzen92
kurzzeitig zu unterbrechen.

6

Das Prinzip des Wettkampfes bringt es mit sich, dass zwei einzelne
Gestalten aus der Masse als deutlich wahrnehmbare Individuen
hervortreten, die als Einzelexistenzen im naturhaften ‚Kampf um's
Dasein` keinen Moment lang Bestand hätte

6

In der dionysischen Natur gleichen sich Schöpfung und Vernichtung in
ihrer untrennbaren Dynamik einander an, die keiner Gestalt Dauer
zugesteht.

9

Genuin kulturelle Faktoren -- Tradition, Bildung und Ethik („Ehre``)
stabilisieren jenen Rhythmus als Maß der Wettkämpfe, der die Natur am
Ausbrechen ihrer vollen Potenz und Zerstörungskraft hindert

9

Denn die Natur kennt keine Feier des Siegers, sie kennt kein Tribunal,
das ihn als solchen markiert und nach einiger Zeit in den nächsten
Zweikampf schickt.

11

Das ästhetische Urteil über Homer als Prototyp des Dichters ist das
Ergebnis eines Jahrhunderte währenden kulturell bedingten Selektions-
und Bildungsprozesses, der nach und nach das Bild Homers als Inbegriff
des Poeten hervortreten lässt

12

Der Held als Einzelner geht schicksalhaft unter; als Heldentypus aber
trotzt er der unabänderlichen dionysischen Wahrheit des Untergangs, da
er mit jeder Aufführung erneut ins Gedächtnis gerufen wird

12

Nietzsches Trennungslinie verläuft nicht cartesianisch zwischen dem
Geist als Unvergänglichem und der Physis als Vergänglichem103, sondern
zwischen anthropologischer apollinischer Gestaltungskraft einerseits und
der dionysischen Natur in ihrer gestaltenfeindlichen Vergänglichkeit
andererseits.

16

Das Instinktive ist insofern nicht als ‚natürlichere` Haltung gegenüber
dem bewussten Handeln zu verstehen, sondern vielmehr als dessen nach und
nach errungene körperliche Selbstverständlichkeit.

17

Die griechische Kultur scheint ein gelungenes Beispiel für einen
Anpassungsprozess an die Zeitlichkeit, ohne durch diese Annäherung die
eigene Identität dem Wandel der Zeit preiszugeben.

\paragraph{annotations: hanna, reents (2016) benn
handbuch.pdf}\label{annotations-hanna-reents-2016-benn-handbuch.pdf}

paper: Hanna and Reents (2016)

\begin{verbatim}
Warning: 'xfun::attr()' is deprecated.
Use 'xfun::attr2()' instead.
See help("Deprecated")
\end{verbatim}

page

note

comment

83

Die Lyrikbände aus den Jahren des Publikationsverbots ab 1938 können
keine breite Wirkung entfalten, da sie nur in wenigen hektographierten
Exemplaren an Freunde wie Oelze und das Ehepaar Leonharda und Julius
Gescher gehen (Biographische Gedichte, 1941) oder als Privatdruck in
sehr kleiner Auflage verbreitet werden (Zweiundzwanzig Gedichte, 1943).

120

Benn schreibt hier Gedankenlyrik, auch wenn sie durch eine »Ambivalenz
von sinnlicher Intensität und abstraktiver Reflexion« gekennzeichnet ist

120

Vier Hauptthemen sind in den Gedichten erkennbar: die elegische
Bestimmung und Beschwörung einer Endzeit; das Poetologische; das
Religiöse; und die Zeitgeschichte, vor allem der Triumph der Kunst über
die Geschichte. Alle vier Themen gehören zusammen, denn die von Benn
beschworene Dichtung ist eine spätzeitliche Dichtung, von Trauer eher
als von Hoffnung bestimmt und erst in der Zeit des Verfalls entstanden.
Sie ist vorwiegend resignativ gefärbt, da der Dichter nicht auf die Welt
wirken kann. (..)(..)Trotz ihrer deutlichen Trennung sind Wirklichkeit
und Kunst begrifflich verwoben.

120

die elegische Bestimmung und Beschwörung einer Endzeit; das
Poetologische; das Religiöse; und die Zeitgeschichte, vor allem der
Triumph der Kunst über die Geschichte

\paragraph{annotations: nietzsche KGW III-2, Nachgelassene Schriften
1870-1873.pdf}\label{annotations-nietzsche-kgw-iii-2-nachgelassene-schriften-1870-1873.pdf}

paper: @

\begin{verbatim}
Warning: 'xfun::attr()' is deprecated.
Use 'xfun::attr2()' instead.
See help("Deprecated")
\end{verbatim}

page

note

comment

3

Nachgelassene Schriften 1870

\paragraph{annotations: Nietzsche, Homers Wettkampf, KGW
III.2.pdf}\label{annotations-nietzsche-homers-wettkampf-kgw-iii.2.pdf}

paper: Nietzsche (1973)

\begin{verbatim}
Warning: 'xfun::attr()' is deprecated.
Use 'xfun::attr2()' instead.
See help("Deprecated")
\end{verbatim}

page

note

comment

1

Der Mensch, in seinen höchsten und edelsten 10 Kräften, ist ganz Natur
und trägt ihren unheimlichen Doppeldiarakter an sidi. Seine furchtbaren
und als unmenschlich geltenden Befähigungen sind vielleicht sogar der
fruchtbare Boden, aus dem allein alle Humanität, in Regungen Thaten und
Werken hervorwachsen kann

2

Warum jauchzte die ganze griechische Welt bei den Kampfbildern der
Ilias?

3

wie sich in Wahrheit vom M o r d e und der Mordsühne aus der Begriff des
griechischen Rechtes entwickelt hat, so nimmt aucái die edlere Kultur
ihren ersten Siegeskranz vom

4

nichts scheidet die S griechische Welt so sehr von der unseren, als die
hieraus abzuleitende F ä r b u n g einzelner ethisciier Begriffe z. B.
der E r i s und des N e i d e s .

5

deren Bedeutung im Gegentheil damit umsdirieben ist, daß mit ihnen der
Mensch η i e den Wettkampf wagen darf, er dessen Seele gegen jedes andre
lebende Wesen eifersüchtig erglüht. Im Kampfe des Thamyris mit den
Musen, 3° des Marsyas mit Apoll, im ergreifenden Sdhicksale der Niobe
erschien das sdireckliche Gegeneinander der zwei Mächte, die nie mit
einander kämpfen dürfen, von Mensch und Gott

marginnote4app://note/B577EF2B-F3C2-4C22-9D87-ACBE417CFEBB

5

deren Bedeutung im Gegentheil damit umsdirieben ist, daß mit ihnen der
Mensch η i e den Wettkampf wagen darf, er dessen Seele gegen jedes andre
lebende Wesen eifersüchtig erglüht. Im Kampfe des Thamyris mit den
Musen, 3° des Marsyas mit Apoll, im ergreifenden Sdhicksale der Niobe
erschien das sdireckliche Gegeneinander der zwei Mächte, die nie mit
einander kämpfen dürfen, von Mensch und Gott

6

O s t r a k i s m o s

scherbengericht

6

O s t r a k i s m o s

scherbengericht

7

Das ist der Kern der hellenischen ``Wettkampf-Vorstellung: sie
verabsciieut die Alleinherrsdiafl: und fürditet ihre Gefahren, sie
begehrt, als S c h u t z m i t t e l gegen das Genie --- ein zweites
Genie

7

Für die Alten aber war das 2 5 Ziel der agonalen Erziehung die
``Wohlfahrt des Ganzen, der staatlichen Gesellsciiaft. Jeder Athener z.
B. sollte sein Selbst im''Wettkampfe soweit entwickeln, als es Athen vom
höchsten Nutzen sei und am wenigsten Schaden bringe.

9

Verhältniß des Wettkampfes zur Conception des Kunstwerkes

10

daß ohne Neid Eifersucht und wettkämpf enden Ehrgeiz der hellenische
Staat wie der hellenische Mensch entartet. Er wird böse und grausam, er
wird rachsüchtig und gottlos, kurz, er wird „vorhomerisch''

10

Sparta und Athen liefern sich an Persien aus, wie es Themistokles und
Alkibiades gethan haben; sie verrathen das Hellenische, nachdem sie den
edelsten hellenischen Grundgedanken, den Wettkampf, aufgegeben haben:
und Alexander, die vergröbernde Copie und Abbreviatur der griechischen
Geschichte, 2j erfindet nun den Allerwelts-Hellenen und den sogenannten
„Hellenismus''

\paragraph{annotations: nietzsche, kga 3-1, geburt d
tragödie.pdf}\label{annotations-nietzsche-kga-3-1-geburt-d-tragodie.pdf}

paper: Nietzsche (2024)

\begin{verbatim}
Warning: 'xfun::attr()' is deprecated.
Use 'xfun::attr2()' instead.
See help("Deprecated")
\end{verbatim}

page

note

comment

7

nietzsche\_kga\_2024

\#init\_Geburt der Tragödie\_\_(..)\#cite\_nietzsche\_kga\_2024\_\_

7

nietzsche\_kga\_2024

\#init\_Geburt der Tragödie\_\_(..)\#cite\_nietzsche\_kga\_2024\_\_

10

Ist Pessimismus n o t h w e n d i g das Zeichen des Niedergangs,
Verfalls, des Missrathenio seins, der ermüdeten und geschwächten
Instinkte? --- wie er es bei den Indern war, wie er es, allem Anschein
nach, bei uns, den „modernen'' Menschen und Europäern ist? Giebt es
einen Pessimismus der S t ä r k e ? Eine intellektuelle Vorneigung für
das Harte, Schauerliche, Böse, Problematische des Daseins aus Wohlsein,
aus 15 überströmender Gesundheit, aus F ü l l e des Daseins? Giebt es
vielleicht ein Leiden an der Ueberfülle selbst?

10

Und die Wissenschaft selbst, unsere Wissenschaft --- ja, was bedeutet
überhaupt, als Symptom des Lebens angesehn, alle Wissenschaft? Wozu,
schlimmer noch, w o h e r --- alle Wissenschaft?

10

Phänomen des Dionysischen? Was, aus ihm geboren, die Tragödie?

10

Was bedeutet, gerade bei den Griechen der besten, stärksten, tapfersten
Zeit, der t r a g i s c h e Mythus?

11

Eine feine Nothwehr gegen --- die W a h r h e i t ?

\#ref\textgreater{} wahrheit und lüge i. a. sinne

11

Eine feine Nothwehr gegen --- die W a h r h e i t ?

\#ref\textgreater{} wahrheit und lüge i. a. sinne

11

das Problem der Wissenschaft kann nicht auf dem Boden der Wissenschaft
erkannt werden

12

ein hochmüthiges und schwärmerisches Buch, das sich gegen das profanum
vulgus der „Gebildeten'' von vornherein noch mehr als gegen das „Volk''
abschliesst,

12

Aufgabe selbst nicht fremder wurde, an welche sich jenes verwegene Buch
zum ersten Male herangewagt hat, --- d i e W i s s e n s c h a f t u n t
e r d e r O p t i k d e s K ü n s t l e r s zu s e h n , d i e K u n s t
a b e r u n t e r d e r d e s L e b e n s

14

Wie? wenn die Griechen, gerade im Reichthum ihrer Jugend, den Willen z u
m Tragischen hatten und Pessimisten waren?

14

das V e r l a n g e n n a c h d e m H ä s s l i c h e n , der gute
strenge Wille des älteren Hellenen zum Pessimismus, zum tragischen
Mythus, zum Bilde alles Furchtbaren, Bösen, Räthselhaften,
Vernichtenden, Verhängnissvollen auf dem Grunde des 5 Daseins, --- woher
müsste dann die Tragödie stammen?

14

welche Bedeutung hat dann, physiologisch gefragt, jener Wahnsinn, aus
dem die tragische wie die komische Kunst erwuchs, der dionysische
Wahnsinn?

15

Was bedeutet, unter der Optik des L e b e n s gesehn, --- die Moral?

15

dass nur als ästhetisches Phänomen das Dasein der Welt g e r e c h t f e
r t i g t ist.

16

Christenthum war von Anfang an, wesentlich und gründlich, Ekel und
Ueberdruss des Lebens am Leben,

16

es giebt zu der rein ästhetischen Weltauslegung und Welt-Rechtfertigung,
wie sie in diesem Buche gelehrt wird, keinen grösseren Gegensatz als die
christliche Lehre, welche n u r moralisch ist und sein will und mit
ihren absoluten Maassen, zum Beispiel schon mit ihrer Wahrhaftigkeit
Gottes, die Kunst, j e d e Kunst in's Reich der L ü g e verweist,

16

die Moral selbst in die Welt der Erscheinung zu setzen, herabzusetzen
und nicht nur unter die „Erscheinungen'' (im Sinne des idealistischen
terminus technicus), sondern unter die „Täuschungen'', als Schein, Wahn,
Irrthum, Ausdeutung, Zurechtmachung, Kunst.

16

das Christenthum als die ausschweifendste Durchfigurirung des
moralischen Thema's, welche die Menschheit bisher anzuhören bekommen
hat.

17

vor der Moral (in Sonderheit christlichen, das heisst unbedingten Moral)
m u s s das Leben beständig und unvermeidlich Unrecht bekommen, weil
Leben etwas essentiell Unmoralisches i s t , --- m u s s endlich das
Leben, erdrückt unter dem Gewichte der Verachtung und des ewigen Nein's,
als begehrens-unwürdig, als unwerth an sich empfunden werden.

17

eine grundsätzliche Gegenlehre und Gegenwerthung des Lebens, eine rein
artistische, eine a n t i c h r i s t l i c h e . Wie sie nennen? Als
Philologe und Mensch der Worte taufte ich sie, nicht ohne einige
Freiheit --- denn wer wüsste den rechten Namen des Antichrist? --- auf
den Namen eines griechischen Gottes: ich hiess sie die d i o n y s i s c
h e .

18

wie müsste eine Musik beschaffen sein, welche nicht mehr romantischen
Ursprungs wäre, gleich der deutschen, --- sondern d i o n y s i s c h e
n ?

18

jetzigen d e u t s c h e n M u s i k , als welche Romantik durch und
durch ist und die ungriechischeste aller möglichen Kunstformen

18

zu einer Zeit, wo der deutsche Geist, der nicht vor Langem noch den
Willen zur Herrschaft über Europa, die Kraft zur Führung Europa's gehabt
hatte, eben letztwillig und endgültig a b d a n k t e und, unter dem
pomphaften Vorwande einer Reichs-Begründung, seinen Uebergang zur
Vermittelmässigung, zur Demokratie und den „modernen Ideen'' machte!

20

Das Lachen sprach ich heilig: ihr höheren Menschen, l e r n t mir ---
lachen!

20

Aber es ist sehr wahrscheinlich, dass es so e n d e t , dass i h r so
endet, nämlich „getröstet'', wie geschrieben steht, trotz aller
Selbsterziehung zum Ernst und zum Schrecken, „metaphysisch getröstet'',
kurz, wie Romantiker enden, c h r i s t l i c h

24

Diesen Ernsthaften diene zur Belehrung, dass ich von der Kunst als der
höchsten Auf15 gäbe und der eigentlich metaphysischen Thätigkeit dieses
Lebens im Sinne des Mannes überzeugt bin, dem ich hier, als meinem
erhabenen Vorkämpfer auf dieser Bahn, diese Schrift gewidmet haben will.

25

Fortentwickelung der Kunst an die Duplicität des A p o l l i n i s c h e
n und des D i o n y s i s c h e n gebunden ist: in ähnlicher Weise, wie
die Generation von der Zweiheit der Geschlechter, bei fortwährendem
Kampfe und nur periodisch eintretender Versöhnung, abhängt.

25

Erkenntniss, dass in der griechischen Welt ein ungeheurer Gegensatz,
nach Ursprung und Zielen, zwischen der Kunst des Bildners, der
apollinischen, und der unbildlichen Kunst der Musik, als der des
Dionysus, besteht: beide so verschiedne Triebe gehen neben einander her,
zumeist im offnen Zwiespalt mit einander und sich gegenseitig zu immer
neuen kräftigeren Geburten reizend, um in ihnen den Kampf jenes
Gegensatzes zu perpetuiren, den das gemeinsame Wort „Kunst'' nur
scheinbar überbrückt;

26

Der schöne Schein der Traumwelten, in deren Erzeugung jeder Mensch
voller Künstler ist, ist die Voraussetzung aller bildenden Kunst, ja
auch, wie wir sehen werden, einer wichtigen Hälfte der Poesie.

26

in dieser Paarung zuletzt das ebenso dionysische als apollinische
Kunstwerk der attischen Tragödie erzeugen.

26

denken wir sie uns zunächst als die getrennten Kunstwelten des T r a u m
e s und des R a u s c h e s ;

27

beherrscht auch den schönen Schein der inneren Phantasie-Welt

27

freudige Nothwendigkeit der Traumerfahrung ist gleichfalls von den
Griechen in ihrem Apollo ausgedrückt worden: Apollo, als der Gott aller
bildnerischen Kräfte, ist zugleich der wahrsagende Gott

27

alle Dinge als blosse Phantome oder Traumbilder vorkommen, als das
Kennzeichen philosophischer Befähigung

27

freudige

27

Wie nun der Philosoph zur Wirklichkeit des Daseins, so verhält sich der
künstlerisch erregbare Mensch zur Wirklichkeit des Traumes;

28

Wenn wir zu diesem Grausen die wonnevolle Verzückung hinzunehmen, die
bei demselben Zerbrechen des principii individuationis aus dem innersten
Grunde des Menschen, ja der Natur emporsteigt, so thun wir einen Blick
in das Wesen des D i o n y s i s c h e n , das uns am nächsten noch
durch die Analogie des R a u s c h e s gebracht wird.

28

darf nicht im Bilde des Apollo fehlen: jene maassvolle Begrenzung, jene
Freiheit von den wilderen Regungen, jene weisheitsvolle Ruhe des
Bildnergottes

28

Apollo als das herrliche Götterbild des principii individuationis
bezeichnen, aus dessen Gebärden und Blicken die ganze Lust und Weisheit
des „Scheines'', sammt seiner Schönheit, zu uns spräche

29

bei dem gewaltigen, die ganze Natur lustvoll durchdringenden Nahen des
Frühlings erwachen jene dionysischen Regungen, in deren Steigerung das
Subjective zu völliger Selbstvergessenheit hinschwindet

30

Diesen unmittelbaren Kunstzuständen der Natur gegenüber ist jeder
Künstler „Nachahmer'', und zwar entweder apollinischer Traumkünstler
oder dionysischer Rauschkünstler oder endlich --- wie beispielsweise in
der griechischen Tragödie --- zugleich Rausch- und Traumkünstler

30

die Kunstgewalt der ganzen Natur, zur höchsten Wonnebefriedigung des
Ur-Einen, offenbart sich hier unter den Schauern des Rausches

32

beschränkte sich das Wirken des delphischen Gottes darauf, dem
gewaltigen Gegner durch eine zur rechten Zeit abgeschlossene Versöhnung
die vernichtenden Waffen aus der Hand zu nehmen.

32

erkennen wir jetzt, im Vergleiche mit jenen babylonischen Sakäen und
ihrem Rückschritte des Menschen zum Tiger und Affen, in den dionysischen
Orgien der Griechen die Bedeutung von Welterlösungsfesten und
Verklärungstagen

32

Es ist die dorische Kunst, in der sich jene majestätisch-ablehnende
Haltung des Apollo verewigt hat.

32

Fast überall lag das Centrum dieser Feste in einer überschwänglichen
geschlechtlichen Zuchtlosigkeit, deren Wellen über jedes Familienthum
und dessen ehrwürdige Satzungen hinweg flutheten; gerade die wildesten
Bestien der Natur wurden hier entfesselt, bis zu jener abscheulichen
Mischung von Wollust und Grausamkeit, die mir immer als der eigentliche
„Hexentrank'' erschienen ist

32

Es war die Versöhnung zweier Gegner, mit scharfer Bestimmung ihrer von
jetzt ab einzuhaltenden Grenzlinien und mit periodischer Uebersendung
von Ehrengeschenken; im Grunde war die Kluft nicht überbrückt.

33

Die Musik des Apollo war dorische Architektonik in Tönen, aber in nur
angedeuteten Tönen, wie sie der Kithara zu eigen sind. Behutsam ist
gerade das Element, als unapollinisch, ferngehalten, das den Charakter
der dionysischen Musik und damit der Musik über25 haupt ausmacht, die
erschütternde Gewalt des Tones, der einheitliche Strom des Melos und die
durchaus unvergleichliche Welt der Harmonie.

34

als sich ihm das Grausen beimischte, dass ihm jenes Alles doch
eigentlich so fremd nicht sei, ja dass sein apollinisches Bewusstsein
nur wie ein Schleier diese dionysische Welt vor ihm verdecke.

36

Wie anders hätte jenes so reizbar empfindende, so ungestüm begehrende,
zum L e i d e n so einzig befähigte Volk das Dasein ertragen können,
wenn ihm nicht dasselbe, von einer höheren Glorie umflossen, in seinen
Göttern gezeigt worden wäre.

36

So rechtfertigen die Götter das Menschenleben, indem sie es selbst leben
---die allein genügende Theodicee!

37

In den Griechen wollte der „Wille'' sich selbst, in der Verklärung des
Genius und der Kunstwelt, anschauen; um sich zu verherrlichen, mussten
seine Geschöpfe sich selbst als verherrlichenswerth empfinden,

37

Die homerische „Naivetät'' ist nur als der vollkommene Sieg der
apollinischen Illusion zu begreifen: es ist dies eine solche Illusion,
wie sie die Natur, zur Erreichung ihrer Absichten, so häufig verwendet

38

Je mehr ich nämlich in der Natur jene allgewaltigen Kunsttriebe und in
ihnen eine inbrünstige Sehnsucht zum Schein, zum Erlöstwerden durch den
Schein gewahr werde, um so mehr fühle ich mich zu der metaphysischen
Annahme gedrängt, dass das WahrhaftSeiende und Ur-Eine, als das ewig
Leidende und Widerspruchsvolle, zugleich die entzückende Vision, den
lustvollen Schein, zu seiner steten Erlösung braucht:

39

muss uns jetzt der Traum als der S c h e i n des S c h e i n s , somit
als eine noch höhere Befriedigung der Urbegierde nach dem Schein hin
gelten

39

die ganze Welt der Qual nöthig ist, damit durch sie der Einzelne zur
Erzeugung der erlösenden Vision gedrängt werde

40

neben der ästhetischen Nothwendigkeit der Schönheit die Forderung des
„Erkenne dich selbst'' und des „Nicht zu viel!'' her, während
Selbstüberhebung und Uebermaass als die eigentlich feindseligen Dämonen
der nicht-apollinischen Sphäre, daher als Eigenschaften der
vor-apollinischen Zeit, des Titanenzeitalters, und der
ausser-apollinischen Welt d.~h. der Barbarenwelt, erachtet wurden.

40

Apollo, als ethische Gottheit, fordert von den Seinigen das Maass und,
um es einhalten zu können, Selbsterkenntniss.

41

denken wir uns, was diesem dämonischen Volksgesange gegenüber der
psalmodirende Künstler des Apollo, mit dem gespensterhaften
Harfenklange, bedeuten konnte! Die Musen der Künste des „Scheins''
verblassten vor einer Kunst, die in ihrem Rausche die Wahrheit sprach,
die Weisheit des Silen rief Wehe! Wehe! aus gegen die heiteren Olympier

41

nur in einem unausgesetzten Widerstreben gegen das titanisch-barbarische
Wesen des Dionysischen konnte eine so trotzig-spröde, mit Bollwerken
umschlossene Kunst, eine so kriegsgemässe und herbe Erziehung, ein so
grausames und rücksichtsloses Staatswesen von längerer Dauer sein

42

Homer, der in sich versunkene greise Träumer, der Typus des
apollinischen, naiven Künstlers, sieht nun staunend den
leidenschaftlichen Kopf des wild durch's Dasein getriebenen 3°
kriegerischen Musendieners Archilochus

42

das erhabene und hochgepriesene Kunstwerk der a t t i s c h e n T r a g
ö d i e und des dramatischen Dithyrambus, als das gemeinsame Ziel beider
Triebe, deren io geheimnissvolles Ehebündniss

42

Wenn auf diese Weise die ältere hellenische Geschichte, im Kampf jener
zwei feindseligen Principien, in vier grosse Kunststufen zerfällt: so
sind wir jetzt gedrängt, weiter nach dem letzten Plane dieses Werdens
und Treibens zu fragen, falls uns S nicht etwa die letzterreichte
Periode, die der dorischen Kunst, als die Spitze und Absicht jener
Kunsttriebe gelten sollte: und hier bietet sich unseren Blicken das
erhabene und hochgepriesene Kunstwerk der a t t i s c h e n T r a g ö d
i e und des dramatischen Dithyrambus, als das gemeinsame Ziel beider
Triebe, deren io geheimnissvolles Ehebündniss, nach langem
vorhergehenden Kampfe, sich in einem solchen Kinde --- das zugleich
Antigone und Kassandra ist --- verherrlicht hat.

42

dem „objectiven'' Künstler der erste „subjective'' entgegen gestellt sei

43

in jeder Art und Höhe der Kunst vor allem und zuerst Besiegung des
Subjectiven, Erlösung vom „Ich'' und Stillschweigen jedes individuellen
Willens und Gelüstens fordern, ja ohne Objectivität, ohne reines
interesseloses Anschauen nie an die geringste wahrhaft künstlerische
Erzeugung glauben können.

44

jetzt aber wird diese Musik ihm wieder wie in einem g l e i c h n i s s
a r t i g en T r a u m b i l d e , unter der apollinischen
Traumeinwirkung sichtbar. Jener bild- und begrifflose Wiederschein des
Urschmerzes in der Musik, mit seiner Erlösung im Scheine, erzeugt jetzt
eine zweite Spiegelung, als einzelnes Gleichniss oder Exempel.

44

Der Plastiker und zugleich der ihm verwandte Epiker ist in das reine
Anschauen der Bilder versunken. Der dionysische Musiker ist ohne jedes
Bild völlig nur selbst Urschmerz und Urwiederklang desselben

47

dass der ganze Gegensatz, nach dem wie nach einem Werthmesser auch noch
Schopenhauer die Künste eintheilt, der des Subjectiven und des
Objectiven, überhaupt in der Aesthetik ungehörig ist, da das Subject,
das wollende und seine egoistischen Zwecke fördernde Individuum nur als
Gegner, nicht als Ursprung der Kunst gedacht werden kann.

47

die Lyrik als eine unvollkommen erreichte, gleichsam im Sprunge und
selten zum Ziele kommende Kunst charakterisirt wird, ja als eine
Halbkunst, deren W e s e n darin bestehen solle, dass das Wollen und das
reine Anschauen d.~h. der unaesthetische und der aesthetische Zustand
wundersam durch einander gemischt seien?

47

denn nur als a e s t h e t i s c h e s P h ä n o m e n ist das Dasein
und die Welt ewig g e r e c h t f e r t i g t :

109

Wir sind wirklich in kurio zen Augenblicken das Urwesen selbst und
fühlen dessen unbändige Daseinsgier und Daseinslust; der Kampf, die
Qual, die Vernichtung der Erscheinungen dünkt uns jetzt wie nothwendig,
bei dem Uebermaass von unzähligen, sich in's Leben drängenden und
stossenden Daseinsformen, bei der überschwänglichen Frucht15 barkeit des
Weltwillens;

\paragraph{annotations: nietzsche, lyrik, benn, benne.zittel
(2017).pdf}\label{annotations-nietzsche-lyrik-benn-benne.zittel-2017.pdf}

paper: Hertel (2017)

\begin{verbatim}
Warning: 'xfun::attr()' is deprecated.
Use 'xfun::attr2()' instead.
See help("Deprecated")
\end{verbatim}

page

note

comment

507

1907 kamen Benn in der von Hans Benzmann zusammengestellten Anthologie
Moderne deutsche Lyrik Texte von 165 Autoren vor Augen -- und er
erinnert sich: »Else Lasker-Schüler war schon vertreten, der junge
Rilke, der junge Hofmannsthal, auch Nietzsche war dabei«

508

Buch von Ernst Bertram -- ›Nietzsche, Versuch einer Mythologie‹,
erschienen 1918« (SW V, 199).3 Unter den Gedichten, die in diesem Buch
von Bertram behandelt wurden, fand er u. a. »Sils-Maria« und »Das Wort«.
In dem letztgenannten Gedicht diagnostiziert Nietzsche, dass das
dichterische Wort »oft schon an bösen Blicken« stirbt

508

Benns Affinität zu Nietzsche wuchs scheinbar in den Phasen als die
»kultische Verehrung im ganzen« nachließ »und überhaupt die Fieberkurve
der Nietzsche-Rezeption deutlich« sank

508

Lasker-Schüler war eine emphatische Anhängerin von Nietzsche. 1907
konnte Benn z. B. in ihrem Gedicht »Weltende« lesen: »als ob der liebe
Gott gestorben war«

508

Rückblickend stellte sie fest: »Nietzsche hat die Sprache geschaffen, in
der wir alle dichten«

508

Im Kontext der damaligen Sprachkrise5 stellte der Kontrast zwischen
fadem Alltagsgerede einerseits sowie andererseits einer formbewussten,
artistischen Sprache, welche lebendige Worte heranspringen lässt, eine
Herausforderung dar. Benn war empfänglich für an- und aufregende Worte,
die Blutdruck und Neurotransmitter in Wallung bringen

508

Seine Rezeption war u. a. durch Empathie mit dem unzeitgemäßen
Einzelgänger verbunden und wurde von Reflexionen begleitet, die
Kunst-Theorie und Praxiserfahrungen abglichen

509

Klabund hatte bereits 1912 die Fröhliche Wissenschaft gelesen. Er
schrieb dann z. B. 1920: »Nietzsche {[}\ldots{]} ist mit der
musikalischen und rhythmischen Prosa seines ›Zarathustra‹ der
Lehrmeister der jungen und jüngsten Dichtung geworden.

509

Gedicht »Vereinsamt« (bzw. »Die Krähen schrei'n«). (..)(..)Ein Vers aus
diesem Gedicht hatte Benn jahrelang beschäftigt. Er erklärte: »Ich
schüttle mir das nicht aus dem Ärmel, was ich jetzt sage, ich habe
jahrelang darüber nachgedacht, ich habe jahrelang über den Vers
nachgedacht: ›Wer das verlor, was ich verlor, macht nirgends halt‹«

509

beide waren fasziniert von Nietzsche, von seinem »Stil!«, von den
»tödlich beschwingten Rhythmen und Akzenten

510

Ich glaube, daß Lyrik ohne Theorie der Lyrik überhaupt bei einem
größeren Lyriker nicht mehr möglich ist«

510

Benn suchte Zeichen-Material, ließ sich ansprechen (von sprachlichen
Reizen verführen) und hatte insbesondere bestimmte Worte im Visier, die
er zum Gedicht-Bau nutzen konnte. Der Rezeptionsmodus bildete also
zugleich das Paradigma für seine Lyrik-Produktion. Gedankliche
Vorarbeiten, Notizen, Essays gehören zum modernen Dichten. Benn berief
sich mit dieser Auffassung ausdrücklich auf Nietzsche und erläuterte,
dass er »überhaupt die Kombination von Lyrik und Essay« als »das
spezifisch Moderne« empfinde

510

In seinem Gedicht »NUR ZWEI DINGE« (1953) finden sich Verse, die auf
jene »ewige Frage: wozu?« (SW I, 320) reagieren. (..)(..)Hier könnte
also ein weiterer Bezug zu dem Nietzsche-Gedicht »Vereinsamt«
hergestellt werden

510

wurde von Benn Gott als (vollendete) Form11 bestimmt und Nietzsches
These ernst genommen, der zufolge Dasein und Welt nur »als aesthetisches
Phänomen« gerechtfertigt sind12

514

Bei den Deutungsversuchen sollte jedoch nicht unbeachtet bleiben, welche
Bedeutungen Nietzsche und Benn dem Leid bzw. dem Leiden beigemessen
haben. Nietzsche verwies schon in der Geburt der Tragödie darauf, dass
das »Talent zum Leiden und zur Weisheit des Leidens« als ein Korrelat
zum künstlerischen Talent angesehen werden kann

515

Auch Benn konnte durch den Sprach-Rhythmus gewichtige Themen (Zeit,
Schicksal, Tod, Leiden) mit geistiger Leichtigkeit ansprechen. Dabei
gilt freilich für ihn (wie für Hölderlin), dass seine Verse »substanzlos
sind, nahezu ein Nichts, um ein Geheimnis geschmiedet, das nie
ausgesprochen wird und das sich nie enthüllt« (SW III, 329). Benn
serviert keine Botschaften, die eindeutig dechiffrierbar sind, sondern
bringt Semantik zum Schwingen

518

drückte er u. a. in folgendem Vers aus: »man sage nicht, der Geist kann
es erreichen, / er gibt nur manchmal kurzbelichtet Zeichen« (SW I, 285).
Diese Sprach-Zeichen sind nicht auf provokante, aufregende
›Wallungsworte‹ beschränkt. Ihre Semantik berührt weiterreichende
Dimensionen30. Im 20. Jahrhundert gab es nämlich neue Möglichkeiten für
Diskurse der Dichtkunst. Benn konnte in dem Freiraum, den Nietzsche
(mit) erschlossen hatte, weit ausholen. Die »Befreiung von Zwecken« ist
»zum letzten Zweck der Lyrik in der Moderne« geworden

519

Nietzsche und Benn nutzten in ihren Gedichten diese erweiterten
ästhetischen Möglichkeiten ausgiebig. (..)(..)Laut Schlaffer schafft der
Dichter »eine moderne Entsprechung zur archaischen Geistersprache: eine
geheimnisvolle Kommunikation mit einem ihm zugänglichen, allen anderen
verborgenen Adressaten. Dieser Adressat ist niemand anders als der
Dichter selbst. Er, der letzte Repräsentant jenseitiger Geister in einer
entgeisterten Welt, führt im Gedicht ein Selbstgespräch«

519

Außerdem ist für Benn »das moderne Gedicht, das absolute Gedicht
{[}\ldots{]} an niemanden gerichtet, ein Gedicht aus Worten, die Sie
faszinierend montieren« (SW VI, 239). Diese »monologische« Dichtkunst
ist »an die Muse gerichtet, und diese ist unter anderem dazu da, die
Tatsache zu verschleiern, daß Gedichte an niemanden gerichtet sind«
(ebd.). War Nietzsche für Benn eine ›Muse‹? Man kann auch fragen:
Betreibt Benn mit seinen lyrischen Anschlüssen an (und gedanklichen
Antworten auf) Nietzsche eine Art Ahnenkult? Der Sprach-Baumeister
Nietzsche hatte selbst ein Muster vorgezeichnet, welches dies nahelegt:
»Für wen baut der Baumeister? {[}\ldots{]} Ich glaube er baut für den
nächsten grossen Baumeister

520

In der Nachfolge von Nietzsche forschte er mit seinen dichterischen
Experimenten insbesondere danach, wie Kreativität generiert wird bzw.
wie die ›schöpferische Lust‹ zur Geltung kommt. Bei seinen Reflexionen
gelangte er an die Grenzen dessen, was sich mit (Natur-)Wissenschaft
beschreiben und erklären lässt. Er versuchte mit magischen Worten und
faszinierenden Formen über diese Barrieren hinauszugelangen. Benn hat
poetische Experimente durchgeführt, um dem (normalerweise) Unsagbaren
nachzuspüren. Dabei schuf er ›merk-würdige‹ Gedichte, die noch heute
Anregungspotenzial besitzen und die quasi als Aufputschmittel für
Schöpfertum wirken können

\paragraph{annotations: roth, roesler (2020) neudeutsche
schule.pdf}\label{annotations-roth-roesler-2020-neudeutsche-schule.pdf}

paper: Von Roth and Roesler (2020)

\begin{verbatim}
Warning: 'xfun::attr()' is deprecated.
Use 'xfun::attr2()' instead.
See help("Deprecated")
\end{verbatim}

page

note

comment

3

Die Neudeutsche Schule -- Phänomen und Geschichte

\#cite\_von\_roth\_neudeutsche\_2020\_\_

3

Die Neudeutsche Schule -- Phänomen und Geschichte

\#cite\_von\_roth\_neudeutsche\_2020\_\_

554

Vgl. Wagner 1849 Kunst und Revolution; Wagner 1850 Kunstwerk der
Zukunft. Als Ideal beschreibt Wagner darin, wie im antiken Griechenland
die Menschen im Theater zusammenkamen, um die Tragödie im Rahmen der
großen Dionysien zu zelebrieren:

2002

Wagner 1849 Kunst und Revolution (..) Wagner, Richard: Die Kunst und die
Revolution, Leipzig 1849, in: Wagner-Schriften 3, S. 8--41.

\paragraph{annotations: wagner (1849), kunst und
revolution.pdf}\label{annotations-wagner-1849-kunst-und-revolution.pdf}

paper: @

\begin{verbatim}
Warning: 'xfun::attr()' is deprecated.
Use 'xfun::attr2()' instead.
See help("Deprecated")
\end{verbatim}

page

note

comment

1

Wagner, Richard, Die Kunst und die Revolution

\#init\_Wagner, Die Kunst und die Revolution\_

1

Wagner, Richard, Die Kunst und die Revolution

\#init\_Wagner, Die Kunst und die Revolution\_

12

Wir fönnen bei einigem Nac)denfen in unferer Runft feinen Gdritt tbun,
obne auf den Jufammenbang Derfelben mit ber Stunft Der Ürieden zu
treffen. Jn Wabrbeit ift unfere moderne Sunft nur ein Olied in der
Stette Der Runftentwicelung Des gefammten Europa, und diefe nimmt ifren
Quêgang bon ben Grieden

\protect\phantomsection\label{refs}
\begin{CSLReferences}{1}{1}
\bibitem[\citeproctext]{ref-benn_gesammelte_1960}
Benn, Gottfried. 1960. \emph{Gesammelte {Werke}. 3, {Gedichte}}. Edited
by Wellershoff. Vol. 3. Gesammelte {Werke}. Limes-Verl.

\bibitem[\citeproctext]{ref-benn_lyrik_1971}
Benn, Gottfried. 1971. \emph{Lyrik Und {Prosa}, {Briefe} Und
{Dokumente}: Eine {Auswahl}}. 8., erw. Aufl. Limes-{Paperback}.
Limes-Verl.

\bibitem[\citeproctext]{ref-benn_gesammelte_1989}
Benn, Gottfried. 1989. \emph{Gesammelte {Werke} in Der {Fassung} Der
{Erstdrucke}. {[}3{]}, {Essays} Und {Reden} in Der {Fassung} Der
{Erstdrucke}}. 1. Auflage. Edited by Hillebrandt. Fischer
{Taschenbücher} 5233.

\bibitem[\citeproctext]{ref-gunther_rhythmus_2008}
Günther, Friederike Felicitas. 2008. {``Rhythmus Beim Frühen
{Nietzsche}.''} In \emph{Rhythmus Beim Frühen {Nietzsche}}. De Gruyter.
\url{https://doi.org/10.1515/9783110210149}.

\bibitem[\citeproctext]{ref-hanna_benn-handbuch_2016}
Hanna, Christian M., and Friederike Reents. 2016. \emph{Benn-{Handbuch}:
{Leben}, {Werk}, {Wirkung}}. 1. Aufl. 2016 edition. J.B. Metzler Verlag.
\url{https://doi.org/10.1007/978-3-476-05310-7}.

\bibitem[\citeproctext]{ref-hertel_modernes_2017}
Hertel, Thomas. 2017. {``Ein Modernes ›{Geistergespräch}‹: {Gottfried}
{Benns} {Auseinandersetzung} Mit {Gedichten} von {Friedrich}
{Nietzsche}.''} In \emph{Nietzsche Und Die {Lyrik}: {Ein} {Kompendium}},
edited by Christian Benne and Claus Zittel. J.B. Metzler.
\url{https://doi.org/10.1007/978-3-476-05596-5_34}.

\bibitem[\citeproctext]{ref-nietzsche_homers_1973}
Nietzsche, Friedrich. 1973. {``Homers {Wettkampf}.''} In \emph{{KGA}
{III}-2. {Band} 2 {Nachgelassene} {Schriften} 1870 - 1873}, edited by
Colli/Montinari, vol. 2. {KGW}, III.2. De Gruyter.
\url{https://doi.org/doi:10.1515/9783110835304}.

\bibitem[\citeproctext]{ref-nietzsche_kga_2024}
Nietzsche, Friedrich. 2024. \emph{{KGA}. {Nietzsche} {Werke}: {Band} 1,
{Abteilung} 3. {Die} {Geburt} Der {Tragödie}. {Unzeitgemäße}
{Betrachtungen} {I} - {III} (1872 - 1874) {Kritische} {Gesamtausgabe}}.
De Gruyter.

\bibitem[\citeproctext]{ref-von_roth_neudeutsche_2020}
Von Roth, Dominik, and Ulrike Roesler, eds. 2020. \emph{Die
{Neudeutsche} {Schule} -- {Phänomen} Und {Geschichte}: {Quellen} Und
{Kommentare} Zu Einer Zentralen Musikästhetischen {Kontroverse} Des 19.
{Jahrhunderts}}. J.B. Metzler.
\url{https://doi.org/10.1007/978-3-476-04923-0}.

\end{CSLReferences}




\end{document}
